\subsection{Трекинг человека и жестов}

Используя библиотеку \textbf{MediaPipe} от Google, мы можем распознавать позу человека (скелет) в реальном времени даже на обычном процессоре. Напомним: \texttt{mediapipe} в браузерном симуляторе недоступна, поэтому примеры этой главы запускаются только на компьютере с реальным дроном.

\subsubsection*{Как это работает?}
1. Получаем кадр с камеры.
2. MediaPipe находит ключевые точки (суставы): плечи, локти, запястья, нос и т.д.
3. Анализируя взаимное расположение точек, мы можем определить жест.

\subsubsection*{Пример: Управление жестами}
В примере \texttt{human-tracking.py} реализована сложная логика управления. Давайте разберем основные идеи.

\textbf{Определение жеста <<Руки вверх>>:}
Если запястья ($y_{wrist}$) выше носа ($y_{nose}$), значит человек поднял руки.

\begin{codebox}[Логика жеста]
\begin{verbatim}
if (left_wrist.y < nose.y) and (right_wrist.y < nose.y):
    print("Жест: Взлет!")
    drone.takeoff()
\end{verbatim}
\end{codebox}

\subsubsection*{Пропорциональное управление (P-регулятор)}
Чтобы дрон поворачивался за человеком, мы вычисляем ошибку: насколько центр человека смещен от центра кадра.
$$ Error = X_{person} - X_{center} $$
Управляющее воздействие (скорость поворота) пропорционально ошибке:
$$ YawRate = K_p \times Error $$
Где $K_p$ — коэффициент чувствительности.

\begin{taskbox}[Джедай]
Напишите упрощенную программу:
\begin{enumerate}
    \item Дрон стоит на земле (не взлетает).
    \item Камера включена.
    \item Если вы поднимаете левую руку — дрон зажигает левый светодиод красным.
    \item Если правую — правый зеленым.
    \item Если скрещиваете руки (расстояние между запястьями мало) — мигает всеми цветами.
\end{enumerate}
\end{taskbox}
